\documentclass[11pt]{article}

\usepackage[margin=1in]{geometry}
\usepackage{amsmath}
\usepackage{booktabs}
\usepackage[colorlinks=true,linkcolor=blue,urlcolor=blue]{hyperref}

\setlength{\emergencystretch}{2em}
\newcommand{\ngvla}{ngVLA}
\newcommand{\arcsec}{^{\prime\prime}}

\title{How the \ngvla\ Subarray Calculator Constructs a Ranked List}
\author{Short technical note accompanying the subarray-selection memo}
\date{21 August 2026}

\begin{document}
\maketitle

\begin{abstract}
This note describes the query-time construction of the ranked subarray list,
assuming that the robust--taper measurements and subarray metadata already
exist.  A 30~GHz request is used so that the frequency scaling is transparent.
The main point is that the calculator does not rank subarrays by inefficiency
$\eta$ alone.  It first finds one feasible imaging operating point for each
subarray, then ranks those points by integration time or antenna hours.
\end{abstract}

\section{Inputs and conventions}

For each offered subarray $i$, the calculator needs metadata
$(N_i,b_{\min,i},C_{{\rm LAS},i},\theta_{{\rm cap},i})$ and measured lattice
points
\[
  p=(\theta,\eta,s,{\rm weighting},{\rm robust},{\rm taper}).
\]
Here $s$ is the PSF skirt at one fitted FWHM, and
\begin{equation}
  \eta(p)=\frac{\sigma_{\rm dirty}(p)}
  {\sigma_{\rm dirty}({\rm natural},0)}
\end{equation}
is normalized to the natural, untapered rms of the \emph{same} subarray.
Thus equal $\eta$ does not imply equal absolute sensitivity for subarrays with
different antenna counts.

The registered configuration and observing mode select one metadata table and
one curve table.  The default is Rev~F, four-hour track.  The selector evaluates
every metadata row independently; it does not first construct a single envelope
across different subarrays.

At a general observing frequency $f$, the reference-curve coordinate is
\[
  \theta_{30}=\theta_{\rm req}\frac{f}{30~{\rm GHz}}.
\]
For the 30~GHz case in this note, $\theta_{30}=\theta_{\rm req}$ and returned
tapers also need no scaling.  The current model treats $\eta$ and skirt as
frequency-invariant.

\section{One candidate from each subarray}

\subsection{The allowed beam interval}

Let $m$ be the fractional beam tolerance (currently 0.05 by default), and let
$\theta_0=\theta_{\rm req}f/(30~{\rm GHz})$ be the central requested beam in
the 30-GHz reference coordinates of the lattice.  Define
\begin{equation}
 \theta_- = \theta_0(1-m),\qquad
 \theta_+ = \theta_0(1+m),\qquad
 I_\theta=[\theta_-,\theta_+].
 \label{eq:beam-window}
\end{equation}
At 30~GHz, $\theta_0=\theta_{\rm req}$.  The calculator does \emph{not}
first set $\theta=\theta_0$ and then minimize $\eta$.  Beam size is itself an
optimization variable constrained to $I_\theta$, so the returned
$\theta_i^*$ can be the central request, either interval boundary, or an
interior value.

\subsection{Interpolation on a two-corner lattice edge}

First consider two adjacent measured lattice points
\begin{align*}
 p_a&=(\theta_a,\eta_a,s_a,r_a,t_a),&
 p_b&=(\theta_b,\eta_b,s_b,r_b,t_b),
\end{align*}
where $r$ is Briggs robust and $t$ is uv taper.  The points are joined only if
they are adjacent along one lattice coordinate and pass the engine's gap
checks.  For any $\theta$ bracketed by $\theta_a$ and $\theta_b$, put
\begin{equation}
 z(\theta)=
 \frac{\ln\theta-\ln\theta_a}{\ln\theta_b-\ln\theta_a},
 \qquad 0\le z\le1.
 \label{eq:edge-fraction}
\end{equation}
The interpolated sensitivity cost and skirt are
\begin{align}
 \eta(\theta)&=\exp\!\left[(1-z)\ln\eta_a+z\ln\eta_b\right]
               =\eta_a^{1-z}\eta_b^z,\label{eq:edge-eta}\\
 s(\theta)&=(1-z)s_a+zs_b.\label{eq:edge-skirt}
\end{align}
The logarithmic coordinates for $\theta$ and $\eta$ are deliberate: both are
positive quantities that span large dynamic ranges, and their changes across
weighting and taper are more naturally described by ratios (locally,
power-law behavior) than by absolute differences.  Linear interpolation in
their logarithms therefore treats equal fractional changes equally and
preserves positivity; skirt is interpolated linearly because it can pass
through zero or become slightly negative.
On a fixed-robust taper-row edge, $r_a=r_b$ and
\begin{equation}
 r(\theta)=r_a,\qquad
 t(\theta)=
 \begin{cases}
  \exp[(1-z)\ln t_a+z\ln t_b],&t_a>0\ \hbox{and}\ t_b>0,\\
  (1-z)t_a+zt_b,&t_a=0\ \hbox{or}\ t_b=0.
 \end{cases}
 \label{eq:edge-row-parameters}
\end{equation}
The second line is used on an edge touching the untapered column because
$\ln0$ is undefined.  This is the complete meaning of the earlier word
``normally.''  On a fixed-taper Briggs-column edge, $t_a=t_b$ and
\begin{equation}
 t(\theta)=t_a,\qquad r(\theta)=(1-z)r_a+zr_b.
 \label{eq:edge-column-parameters}
\end{equation}
Natural-weighting points participate only in fixed-weighting taper rows;
four-corner cells, described next, are Briggs-only.

\subsection{Interpolation inside a four-corner Briggs cell}

Let a valid lattice cell have robust nodes $r_0<r_1$ and taper nodes
$t_0<t_1$, with measured corners
\begin{align*}
 p_a&=(r_0,t_0;\theta_a,\eta_a,s_a),&
 p_b&=(r_1,t_0;\theta_b,\eta_b,s_b),\\
 p_c&=(r_0,t_1;\theta_c,\eta_c,s_c),&
 p_d&=(r_1,t_1;\theta_d,\eta_d,s_d).
\end{align*}
All four corners must be present and lattice-adjacent; the engine does not
form a cell across an intervening sample, a pruned taper gap, or a disallowed
robust or resolution step.  Introduce normalized coordinates
$q,u\in[0,1]$ along the robust and taper directions, respectively, and the
bilinear weights
\begin{equation}
 w_a=(1-q)(1-u),\quad w_b=q(1-u),\quad
 w_c=(1-q)u,\quad w_d=qu,
 \qquad \sum_j w_j=1.
 \label{eq:cell-weights}
\end{equation}
Writing $x_j=\ln\theta_j$ and $e_j=\ln\eta_j$, the interpolated physical
quantities are
\begin{align}
 x(q,u)&=\sum_{j\in\{a,b,c,d\}}w_jx_j,&
 \theta(q,u)&=\exp x(q,u),\label{eq:cell-theta}\\
 e(q,u)&=\sum_{j\in\{a,b,c,d\}}w_je_j,&
 \eta(q,u)&=\exp e(q,u),\label{eq:cell-eta}\\
 s(q,u)&=\sum_{j\in\{a,b,c,d\}}w_js_j.&&
 \label{eq:cell-skirt}
\end{align}
An imaging cell is usable only when skirt is measured at all four corners.
The corresponding \texttt{tclean} parameters are
\begin{align}
 r(q)&=r_0+q(r_1-r_0),\label{eq:cell-robust}\\
 t(u)&=
 \begin{cases}
  \exp[(1-u)\ln t_0+u\ln t_1],&t_0>0,\\
  (1-u)t_0+ut_1,&t_0=0.
 \end{cases}
 \label{eq:cell-taper}
\end{align}
Thus the cell is bilinear in $\ln\theta$, $\ln\eta$, and $s$ as functions of
$(q,u)$; reported robust is linear in $q$, and reported taper is geometric in
$u$ except in a cell touching the untapered column.  The implementation rounds
the returned robust value to three decimal places.

For a particular trial beam $\theta_{\rm trial}$, the cell points that make exactly that
beam lie on the one-dimensional iso-resolution locus
\begin{equation}
 x(q,u)=x_{\rm trial},\qquad x_{\rm trial}=\ln\theta_{\rm trial}.
 \label{eq:iso-resolution}
\end{equation}
For example, define
\begin{align*}
 A&=x_a,&B&=x_b-x_a,&C&=x_c-x_a,&
 D&=x_a-x_b-x_c+x_d.
\end{align*}
Then $x(q,u)=A+Bq+Cu+Dqu$, and, away from the degenerate denominator,
\begin{equation}
 u(q;x_{\rm trial})=\frac{x_{\rm trial}-A-Bq}{C+Dq}.
 \label{eq:cell-isolocus}
\end{equation}
Only values with $0\le q\le1$ and $0\le u(q;x_{\rm trial})\le1$ are in the cell.
Substitution of equation~\eqref{eq:cell-isolocus} into
equations~\eqref{eq:cell-eta}--\eqref{eq:cell-skirt} gives the corresponding
$\eta(q)$ and $s(q)$ along that fixed-beam locus.  Once the optimizer chooses
$q_*$, equations~\eqref{eq:cell-robust} and~\eqref{eq:cell-taper} give the
returned robust $r_*$ and taper $t_*$, and
$\eta_*=\eta(q_*,u_*)$, $s_*=s(q_*,u_*)$.  The implementation treats the
degenerate case $C+Dq=0$ separately rather than dividing by zero.

\subsection{How the beam interval enters the point-source and imaging choices}

For point-source selection, for a given subarray and trial resolution
$\theta_{\rm trial}$, the corresponding $(r_*,t_*)$ that minimize $\eta$ are
computed over the grid of simulation results, including valid interpolated
edges and four-corner cells.  Within a cell, $\ln\theta$, $\ln\eta$, and $s$
are bilinearly interpolated in the normalized robust and taper coordinates
$(q,u)$, yielding the corresponding $(\eta_*,s_*)$ and robust--taper setting.
If a range of resolutions is allowed, the cell-interior calculation is run at
33 values of $\theta_{\rm trial}$ spaced uniformly in $\ln\theta$ within
$[\theta_-,\theta_+]$, and the candidate that minimizes $\eta$ is selected.
For imaging, whenever possible this $\eta$ minimization is subject to the
additional criterion $s\le s_0$.  If all candidate PSFs have $s>s_0$, then
$s$ is minimized instead of $\eta$ (with $\eta$ breaking skirt ties), the
so-called relaxed mode.

In a four-corner cell, the fixed-resolution calculation just described is the
minimization of $\eta(q,u)$ along the iso-resolution locus
$\ln\theta(q,u)=\ln\theta_{\rm trial}$ derived above.

Let $\mathcal P_i(\theta)$ denote all measured, edge-interpolated, and
cell-interior operating points of subarray $i$ that produce beam $\theta$.
For a point-source request, the selected operating point and its achieved beam
are jointly
\begin{equation}
 (\theta_i^*,p_i^*)\in
 \mathop{\rm arg\,min}_{\substack{\theta\in I_\theta\\
                                  p\in\mathcal P_i(\theta)}}\eta(p).
 \label{eq:point-source-choice}
\end{equation}
For an imaging request with required skirt $s_0$, define the strict feasible
set
\begin{equation}
 \mathcal F_i=\{(\theta,p):\theta\in I_\theta,
               \ p\in\mathcal P_i(\theta),\ s(p)\le s_0\}.
 \label{eq:strict-set}
\end{equation}
If $\mathcal F_i$ is nonempty, the strict answer is
\begin{equation}
 (\theta_i^*,p_i^*)\in
 \mathop{\rm arg\,min}_{(\theta,p)\in\mathcal F_i}\eta(p).
 \label{eq:strict-choice}
\end{equation}
This equation is where $I_\theta$ enters the imaging case: $\theta_i^*$ is
chosen at the same time as $\eta_i^*$, $s_i^*$, $r_i^*$, and $t_i^*$, rather
than being fixed to $\theta_0$ beforehand.

If $\mathcal F_i$ is empty and relaxation is enabled, let
\begin{equation}
 s_{i,\min}=\min_{\substack{\theta\in I_\theta\\
                            p\in\mathcal P_i(\theta)}}s(p).
\end{equation}
The relaxed answer first minimizes skirt and then breaks numerical skirt ties
by $\eta$:
\begin{equation}
 (\theta_i^*,p_i^*)\in
 \mathop{\rm arg\,min}_{\substack{\theta\in I_\theta,
                                  \ p\in\mathcal P_i(\theta)\\
                                  s(p)\le s_{i,\min}+10^{-6}}}\eta(p).
 \label{eq:relaxed-choice}
\end{equation}
It is flagged as relaxed and does not meet $s_0$.  With relaxation disabled,
an empty $\mathcal F_i$ excludes the subarray.

The candidate search is finite.  On each two-corner edge the engine evaluates
measured points inside $I_\theta$, crossings of $\theta_-$ and $\theta_+$,
and, for strict imaging, crossings of $s=s_0$.  These contain every edge
minimum because equations~\eqref{eq:edge-eta} and~\eqref{eq:edge-skirt} are
linear in $\ln\theta$ after taking $\ln\eta$.  In a four-corner cell, for a
fixed $\theta_{\rm trial}$ the engine substitutes equation~\eqref{eq:cell-isolocus} and
evaluates interior stationary points of $\eta$, stationary points of $s$, and
crossings of $s=s_0$; the cell edges are already covered by the edge search.
For a nonzero-width $I_\theta$, the current implementation repeats this
fixed-beam calculation at 33 beams spaced uniformly in $\ln\theta$ from
$\theta_-$ through $\theta_+$.  Thus edge minimization is analytic over the
whole interval, while the additional cell-interior minimization over beam is
a deterministic 32-subinterval scan.  A zero-width request evaluates the one
specified beam exactly.  Finally the engine compares all candidates according
to equations~\eqref{eq:point-source-choice}, \eqref{eq:strict-choice}, or
\eqref{eq:relaxed-choice}.  It retains the winning measured anchors and the
interpolated \texttt{tclean} weighting, robust, and taper values.

No interpolation family extrapolates beyond its measured endpoints, and no
edge or cell bridges a pruned gap.

The preceding candidate-search description is the Python
\texttt{Envelope.solve} path.  The in-page browser calculator uses the same
interpolation results, but it does not rerun the cell solve for each form
submission.  At report-build time the engine evaluates each subarray at
beams spaced by 0.01 dex in $\log_{10}\theta$ (plus coverage endpoints) with
zero beam margin, and embeds those curve points in the page.  At query time
the browser restricts that stored set to $I_\theta$ and applies
equations~\eqref{eq:point-source-choice}, \eqref{eq:strict-choice}, and
\eqref{eq:relaxed-choice} to those points.  Thus the browser's
$\theta_i^*$ is one of the embedded curve beams; its $\eta_i^*$, $s_i^*$,
$r_i^*$, and $t_i^*$ were calculated from the edge or four-corner equations
above when that curve point was built.

Two other feasibility gates apply.  A heterogeneous subarray is excluded when
the entire beam window is coarser than its beam cap; if the interval only
overlaps the cap, it is clipped at the cap.  A requested largest angular scale
is checked using
\begin{equation}
 {\rm LAS}_i=C_{{\rm LAS},i}\frac{\lambda}{b_{\min,i}}206264.8.
\end{equation}
At 30~GHz and $C_{\rm LAS}=0.8$, this is approximately
$1649/b_{\min,i}[{\rm m}]$ arcsec.  LAS is reported but is not a gate when the
request omits it.  There is no query-time $\eta\le5$ rejection.

\section{Sensitivity cost and the role of antenna count}

For a supplied flux-density goal $g$, the sensitivity wrapper obtains the
one-hour natural rms $\sigma_{{\rm nat},i}$ from the Memo~96 calculator and
applies the curve-derived inefficiency externally:
\begin{equation}
 \sigma_i(t)=\eta_i\sigma_{{\rm nat},i}(1~{\rm h})
             \sqrt{\frac{1~{\rm h}}{t}},
 \qquad
 t_i=\left(\frac{\eta_i\sigma_{{\rm nat},i}(1~{\rm h})}{g}\right)^2{\rm h}.
\end{equation}
The resource cost is then
\begin{equation}
 A_i=N_i t_i
\end{equation}
in antenna hours.  A brightness-temperature goal is converted to Jy/beam
using the candidate's achieved beam and request frequency before evaluating
the same expression.

There is no double counting of $N_i$ in these equations.  The upstream natural
rms already contains
\begin{equation}
 \sigma_{{\rm nat},i}\mathrel{\propto}
 \frac{1}{\sqrt{N_i(N_i-1)t}}.
\end{equation}
The wrapper rescales this reference natural sensitivity only when the backing
calculator does not already represent the candidate being evaluated.  There
are two intelligible cases.  First, a named backing entry can carry the antenna
count of an older configuration: for example, the calculator has a
Rev-D-vintage core entry, not a separate Rev~C core entry, so a Rev~C request
must replace the stored count by the Rev~C metadata count.  Second, some
offered subarrays have no named backing entry at all; for these, the
main-array natural sensitivity is used as a same-dish reference and rescaled
to the candidate's antenna count.  Ordinary Rev~F and Rev~D named entries have
matching stored and actual counts, so their correction factor is exactly one.
If both a donor reference and a revision-specific count are involved, the
intermediate count cancels algebraically, leaving one main-to-actual
correction.  Finally, multiplying $t_i$ by $N_i$ converts elapsed time to the
EOP resource unit; it does not alter the thermal noise.

When no numerical sensitivity goal is supplied, the calculator defines the
common fiducial goal
\begin{equation}
 g_{\rm fid}=\sigma_{{\rm nat},\mathrm{MAIN}}(1~{\rm minute})
\end{equation}
at the requested frequency and spectral setup, and evaluates every subarray
with the same absolute time equation above.  Thus natural, untapered MAIN
takes one minute by construction, while other beams and subarrays have finite,
interpretable costs.  This is still not simply a lowest-$\eta$ ordering: the
natural-rms anchor includes $N_i(N_i-1)$, so a larger subarray can rank ahead
despite a larger taper penalty.

\section{Final ordering and current caveats}

For imaging, feasible candidates are first partitioned into strict results
that meet $s_0$ and relaxed results that do not.  Strict results are always
listed first.  Within each group, the default sort is increasing antenna
hours, with fewer antennas as a tie breaker; the alternative is increasing
integration time.  If no feasible subarray meets $s_0$, the calculator emits
a structured warning and retains the relaxed list.  Python callers can
explicitly request a mixed-pool policy when a custom ranking needs to trade
skirt compliance against cost.

Sensitivity failures cannot enter a relative-cost fallback.  Failure to cost
an explicit science goal excludes the affected candidate.  Failure of the
implicit fiducial leaves the independently determined beam/LAS feasibility
visible, but disables the entire recommendation ranking and returns a
structured cost warning.  Excluded subarrays follow the feasible rows with
their reasons.

\section*{Implementation locations}

The relevant implementation is in \path{sa_select/curves.py},
\path{sa_select/envelope.py}, \path{sa_select/selector.py},
\path{sa_select/heuristics.py}, and \path{sa_select/sensitivity.py}.  The
browser implementation is in \path{curvegen/report_template.html}, with its
data payload assembled by \path{curvegen/make_report.py}.

\end{document}
